\subsection{2 Phase Locking (2PL)}
This type of locking is widely used since many decades and is the canonical implementation of concurrency control.
Of course, alternatives exist and the locking implementation has a huge impact on performance and scalability.

Concurrency control is implemented using locking, typically using a locking table.
The locks cover tuples, tables, indexes, blocks and more.

Locks have different semantics that indicate whether they can be shared or are exclusive.

Compatibility matrix:
\begin{tables}{cccc}{ & NL         & S          & X}
              S       & \checkmark & \checkmark & -  \\
              X       & \checkmark & -          & -  \\
\end{tables}

\subsubsection{The protocol}
\begin{enumerate}
    \item Before accessing an object, a transaction must acquire lock
    \item A transaction acquires a lock only once. Lock upgrades are possible
    \item A transaction is blocked if the lock request cannot be granted according to the compatibility matrix
    \item A transaction goes through two phases:
          \begin{enumerate}
              \item Growth: Acquire locks, but never release locks
              \item Shrink: Release locks, but never acquire locks
          \end{enumerate}
    \item At EOT (commit or abort), all locks must be released.
\end{enumerate}

The reason for the two phases is that 2PL prevents histories where a transaction must be aborted due to concurrency control issues.
An operation is thus not allowed unless we are sure it is safe.

Without the two phases, the following can happen:
$L_1[x] w_1[x] U_1[x] L_2[x] w_2[x] U_2[y] L_2[y] w_2[y]$, etc, thus meaning that if $T_1$ now access $y$ or $x$ for read or write, the history becomes non-serializable.

The basic 2PL doesn't say anything about aborting or committing transactions.
The issue is that if we are to undo the changes of $T_1$, we would also undo some of the changes from $T_2$, which would then in turn not be restorable anymore.

\subsubsection{Strict 2PL}
Strict 2PL combines 2PL with strict histories. This means that we enforce 2PL, but we also keep the read and write locks until the transaction commits or aborts.
This results in serializable and strict histories.
All DB engines implement strict 2PL.


\subsubsection{Deadlocks Detection}
In theory, we could use a wait-for graph, which is however expensive.
For example, for $T_1 \rightarrow T_2 \rightarrow T_3 \rightarrow T_4 \rightarrow T_1$, we could abort $T_3$ to resolve the cycle.

In practice, no wait-for graphs are built. Transactions and queries have timers. If the timer expires, the transaction or query is aborted.
Deadlocks are considered rare enough to prefer making mistakes every now and then (aborting a transaction that was fine),
rather than implementing something as expensive as a wait-for graph.

Implications from this are for example that long transactions may block queries and that those queries may abort because the timer expired.

However, this is only common in poorly written applications.
